\section{The Staged Pandora's Box}
\label{sec:setting}

\paragraph{Candidates, ladders and the collectable set.}
There is a fixed pool $\cX=\{x_1,\ldots,x_N\}$ of candidates, the \emph{arms}. Arm $i$ has a fixed finite
ladder of $T_i$ stages, or rungs; after $b$ advances in total it stands at stage $t_{i,b}\in\{0,\ldots,T_i\}$,
and one action, an \emph{advance}, takes it to stage $t=t_{i,b}+1$ at a cost and reveals a
reading $Y$ in a standard Borel space, with a measurable score $\varsigma_{i,t}(Y)$; in our
examples a reading is its own score. The set $\cK_i\subseteq\{1,\ldots,T_i\}$ is fixed before the run. A stage
$t\in\cK_i$ is \emph{collectable}: its score
enters the incumbent. Write $\Yc=\varsigma_{i,t}(Y)$ if the stage reached is collectable and $\Yc=-\infty$
otherwise, the \emph{collectable reading}. The history state is
$S_b=(\cC_b,t_{1:N,b},\tilde y_b,K_b)$. Here $\cC_b$ records the ordered action and observation
history, including noncollectable readings and any fixed starting context; $t_{1:N,b}$ records
the arms' stages; $\tilde y_b$ is the best collectable reading seen or the initial value $q_0$;
and $K_b$ is cumulative cost. The \emph{planning law}
$p_\theta(\cdot\mid S_b,i)$ is the predictive distribution of arm $i$'s next reading under a
specified predictive model whose parameters $\theta$ are fixed before the run and never refitted; it may condition
on the whole context, so an observation on one arm can change the prediction for every arm.
Advancing $i\in\cI(S_b)$, the set of arms with a stage left, costs $c(S_b,i)$ and reveals
$Y_{b+1}\sim p_\theta(\cdot\mid S_b,i)$; the successor $S_{b+1}=F(S_b,i,Y_{b+1})$ adds the
observation to the context, advances $t_i$ by one, sets
$\tilde y_{b+1}=\max\{\tilde y_b,\Yc_{b+1}\}$ and adds the cost to $K_b$. Availability is
arm-local: advancing $i$ never changes another arm's ladder. A policy $\pi$ chooses, measurably
in the observed history, an action in $\{\mathrm{stop}\}\cup\cI(S_b)$; $\tau_\pi$ is its
stopping time. Randomized policies may use auxiliary randomness independent of the innovations
driving the planning kernel. The searcher always retains earlier collectable readings, so collecting at an
intermediate stage does not require terminating the search.

\paragraph{Utility, values and the Bellman equation.}
The objective is $\sup_\pi\E_{\theta,\pi}[\tilde y_{\tau_\pi}-K_{\tau_\pi}\mid S_0]$. It is the
budget-free utility of \citet{vu2026coststop} for fixed-pool freeze-thaw optimization, with a
cost per stage in place of a cost per increment and only collectable readings entering the
score. The stop-now payoff is $h(S_b)=\tilde y_b-K_b$, the value of $\pi$ is
$V_\pi(S_b)=\E_{\theta,\pi}[\tilde y_{\tau_\pi}-K_{\tau_\pi}\mid S_b]$, and
$Q_\pi(S_b,i)=\E_\theta[V_\pi(F(S_b,i,Y_{b+1}))\mid S_b,i]$ is the value of forcing $i$ once and
then following $\pi$. The optimal value $V^\star=\sup_\pi V_\pi$ satisfies
$V^\star(S)=\max\{h(S),\max_{i\in\cI(S)}Q^\star(S,i)\}$, with $\max\varnothing=-\infty$, and the
ladders being finite, a policy that attains the maximum at every state is optimal.
For comparison with continuation-value conventions, define
$U_\pi(S)=V_\pi(S)+K(S)$. With $S'=F(S,i,Y)$,
\[
 Q_\pi(S,i)+K(S)=-c(S,i)+\E_\theta[U_\pi(S')\mid S,i].
\]
Thus $V_\pi$ includes sunk costs, while $U_\pi$ removes them; the cost of the
forced advance is already included in $K(S')$ inside the definition of $Q_\pi$.

\begin{assumption}[Standing assumptions]\label{assm:standing}
The planning law is a fixed measurable kernel, never refitted during a run; all scores and
incumbents lie in fixed bounds $[y_{\min},y_{\max}]$ chosen in advance; the cost $c(S,i)$ is measurable, determined by
$(S,i)$ before the advance, and satisfies $0<c_{\min}\le c(S,i)\le c_{\max}<\infty$; every ladder
is finite. History states carry the Borel sigma-fields generated by the finite arm labels and
their readings; the successor map is measurable.
\end{assumption}

\paragraph{Finite horizon and conventions.}
We write $q(S)=\tilde y(S)$ and $L(S)=\sum_i(T_i-t_i)$ throughout. Every advance lowers $L$
by one, so every admissible policy stops after at most $L(S)$ further advances, and
finite-action backward induction defines Borel Bellman values and an attaining deterministic
policy, choosing the first action in any fixed measurable tie order. The stated index policy
stops on an index--incumbent tie. We use $(-\infty-r)^+=0$ for finite $r$ and
$\max\varnothing=\sup\varnothing=-\infty$. Appendix~\ref{app:notation} collects the remaining
conventions and compares them with the closest index and inspection models.

\paragraph{Per-arm planning laws.}
A planning law is \emph{per-arm} when there are measurable kernels $P_i(dy\mid s_i)$ and
cost functions $c_i(s_i)$ such that $p_\theta(dy\mid S,i)=P_i(dy\mid s_i)$ and
$c(S,i)=c_i(s_i)$. An arm's \emph{own state} $s_i$ is its stage and complete observation
history. Advancing another arm leaves $s_i$, its predictive kernel and its next cost unchanged;
an arm does not evolve while paused. We write $\nu(s_i)=T_i-t_i$ and abbreviate its cost as
$c(s_i)$. No Gaussian, density, or scalar first-order Markov assumption is imposed.

At a fixed state $S$, iterate $P_i$ to obtain the law $\mathsf P_i^{s_i}$ of the entire
remaining path of arm $i$. The common space used in the independent-arm proofs has law
\[
\mathsf P_S=\bigotimes_{i=1}^N\mathsf P_i^{s_i}.
\]
One can sample these paths in advance and reveal their next coordinates as the policy chooses
arms. Autonomy makes this construction agree with the sequential planning kernels. Revealing a
coordinate of one path conditions that path alone; the other remaining paths keep their laws.
Dependence within a path is unrestricted and its future costs are functions of its successive
own states. Theorem~\ref{thm:index} uses this product construction; Section~\ref{sec:robust}
does not assume it.

\paragraph{Single-arm stopping problems.}
Fix a history state $S$, an available arm $i$, and a depth cap $k\ge1$. Advance only $i$,
obtaining states $S^{(m)}$ by iterating the planning kernel and successor map. Let
$\mathcal G_m$ be the filtration of those observations and $\nu_i=T_i-t_i$ the remaining
stages. An admissible arm run is a stopping time $1\le\sigma\le\min\{k,\nu_i\}$ for this
filtration; it pays every advance up to $\sigma$ and may abandon the arm after any reading.
Write $\mathcal R_i^{(k)}(S)$ for these runs, $M_\sigma$ for the best newly obtained
collectable reading ($-\infty$ if none), and $\kappa_\sigma$ for the run's total
additional cost. Earlier collectable readings remain in the incumbent, rather
than in $M_\sigma$. For an outside value $r$, define
\[
\Phi_i^{(k)}(S;r)=\sup_{\sigma\in\mathcal R_i^{(k)}(S)}
\{\E[(M_\sigma-r)^+\mid S]-\E[\kappa_\sigma\mid S]\},\qquad
G_i^{(k)}(S)=\sup\{r:\Phi_i^{(k)}(S;r)>0\}.
\]
The empty supremum for $G$ is $-\infty$. The run law is the original planning law from $S$;
$r$ is a payoff comparison, not a modification of that law. The full index uses $k=\nu_i$;
exhausted arms have index $-\infty$. If a collectable stage remains, let $\mu_i(S)$ be the
number of advances to the next such stage. The first advance is forced in $\Phi$; the option
of stopping immediately is represented by $\Phi^+=\max\{0,\Phi\}$. The auxiliary value $r$
may lie outside the range of observed scores.

\begin{proposition}[Reservation-value properties]\label{prop:index}
Under Assumption~\ref{assm:standing}, $r\mapsto\Phi_i^{(k)}(S;r)$ is nonincreasing and
$1$-Lipschitz. The index is finite exactly when a collectable stage lies within the cap, in
which case
\[
y_{\min}-\mu_i(S)c_{\max}\le G_i^{(k)}(S)\le y_{\max}-c(S,i).
\]
In every case $\Phi_i^{(k)}(S;r)>0$ if and only if $r<G_i^{(k)}(S)$, and $G_i^{(k)}(S)$ is
nondecreasing in $k$. When the index is finite, $\Phi_i^{(k)}(S;G_i^{(k)}(S))=0$.
\end{proposition}
Appendix~\ref{app:proof-details} gives the proof.

The index is measurable: it is also the supremum of the rational $a$ for which
$\Phi_i^{(k)}(S;a)>0$. Thus selecting a smallest-numbered index maximizer defines a measurable
policy even for continuous or atomic reading laws.

These definitions also apply under a shared planning law: the current history fixes the
starting conditional law, and the hypothetical run observes only the chosen arm. The
independent-arm theorem requires the additional autonomy assumption given above.

